






\section{Proof of Lemmas for Theorem~\ref{thm:alpha}}
\label{sec:proof_alpha}
\subsection{Proof of Lemma~\ref{lem:Vt_norm_upper}}


%\begin{proof}
\proof{}
	Let $\theta = t x $ where $x$ is any unit vector in $\RR^d$ and $t\in \RR^{+}  $ is a scalar. Consider the equation that characterizes the intersection  $\{t x : t\in\RR^{+} \}\cap \cE_t $, namely
	$
	(t x - \hparam_t)^\top V_t (t x - \hparam_t) \leq  \beta_t^2
	$.
	Equivalently, we have
	$
	t^2\|x\|_{V_t}^2 - 2 t  x^\top V_t \hparam_t + \phi_t^2 \leq 0
	$. This quadratic inequality of $t$ has at least one solution if the discriminant is non-negative, i.e.
	$
	4(x^\top V_t \hparam)^2 \geq 4 \|x\|_{V_t}^2 \phi_t^2 
	$.
	Then by direct computation,
	\$
	\|x\|_{V_t} \leq \frac{\sqrt{(x^\top V_t \hparam_t)^2}}{\phi_t} \leq \frac{\sqrt{x^\top x}\sqrt{(\hparam_t)^\top V_t^\top V_t \hparam_t }}{\phi_t} =  \frac{ \|\hparam_t\|_{V_t^2}} {\phi_t	}.
	\$

	Note that $\theta \in \cE_t$ if and only if $\theta$ is on a ray starting from the origin and intersects $\cE_t$ at one or more points. Namely, 
	$\theta = t x $ for some $x$ that satisfies $4(x^\top V_t \hparam)^2 \geq 4 \|x\|_{V_t}^2 \phi_t^2$, we conclude the proof.
	

	
%\end{proof}
\endproof
\subsection{Proof of Lemma~\ref{lem:Vt_norm_lower}}
%\begin{proof}
\proof{}
Note that for any $\theta \in \cH_t \subset\cE_t$, it holds that  $\|\theta\| \leq \|\hparam_t\| +  \beta_t/\lambda_d(V_t)  $. Also, by the construction of $\cH_t$, we have $\|\theta\|_{V_t}  =  \phi_t$. Then by direct computation, we have
	\$
	\|x^\star(\theta )\|_{V_t} = \frac{\|\theta\|_{V_t}}{\|\theta\| } \geq \frac{\phi_t}{\|\hat{\theta}_t\| +  \beta_t/\lambda_d(V_t) } .
	\$

To prove the same result for any $\theta \in \cE_t$, we only need to show there exists $\theta^\prime \in \cH_t$ such that $x^\star(\theta) = x^\star(\theta^\prime)$. To see this,  let $x = \theta/\|\theta\|$ and consider the  intersection  $\{t x : t\in\RR^{+}  \}\cap \cE_t$, which is non-empty by our choice of $\theta$. Similar to the proof of Lemma \ref{lem:Vt_norm_upper}, the discriminant is non-negative, i.e.
	$
	4(x^\top V_t \hparam)^2 \geq 4 \|x\|_{V_t}^2 \phi_t^2 
	$. 

Now consider the intersection  $\{t x : t\in\RR \}\cap \partial\cE_t$, where we let $\partial \cE_t \coloneqq \left \{\theta \in \RR^d: \|\theta - \hparam_t\|_{V_t}= \beta_t(\delta)\right \}$ be the border of the ellisoid. The intersection points are characterized by the solution to 
\#\label{eq:border_intersection}
t^2\|x\|_{V_t}^2 - 2 t  x^\top V_t \hparam_t + \phi_t^2 = 0.
\#


If $
	4(x^\top V_t \hparam)^2 = 4 \|x\|_{V_t}^2 \phi_t^2 
	$ and there is only one intersection point, namely $\theta$ itself, we have
	\$
	0 = t^2\|x\|_{V_t}^2 - 2 t  x^\top V_t \hparam_t + \phi_t^2 = t^2\|x\|_{V_t}^2 - 2 t  \|x\|_{V_t} \phi_t  + \phi_t^2 = (\|tx\|_{V_t} - \phi_t)^2.
	\$
Therefore, we have $\|\theta\|_{V_t} = \|tx\|_{V_t} = \phi_t$, i.e.  $\theta \in \cH_t$. 
	



If $4(x^\top V_t \hparam)^2 > 4 \|x\|_{V_t}^2 \phi_t^2$, it follows that $x^\top V_t \hparam > \|x\|_{V_t}\phi_t$. This inference is valid given that $x^\top V_t \hparam > 0$, which can be verified using Equation \eqref{eq:border_intersection} and noting that $t>0$.
Consider the solutions to \eqref{eq:border_intersection}
\$
t_1 = \frac{x^\top V_t \hparam_t - \sqrt{(x^\top V_t \hparam)^2 - \|x\|_{V_t}^2 \phi_t^2 }  }{\|x\|_{V_t}^2}, \quad t_2 = \frac{x^\top V_t \hparam_t + \sqrt{(x^\top V_t \hparam)^2 - \|x\|_{V_t}^2 \phi_t^2 }  }{\|x\|_{V_t}^2}.
\$
We only need to show $ \|t_1 x\|_{V_t}< \phi_t  < \|t_2 x\|_{V_t}$, then by the continuity of $\|\cdot\|_{V_t}$ and the convexity of $\cE_t$, there exists $t^\prime \in (t_1,t_2)$ such that $\|t^\prime x\|_{V_t} = \phi_t $. Then $\theta^\prime \coloneqq t^\prime x \in \cH_t$ is the desired point. By direct computation, 
\$
\|t_1 x\|_{V_t} = \frac{x^\top V_t \hparam_t - \sqrt{(x^\top V_t \hparam)^2 - \|x\|_{V_t}^2 \phi_t^2 }  }{\|x\|_{V_t}}.
\$
We only need to prove 
\$
x^\top V_t \hparam_t -\|x\|_{V_t}  \phi_t  \leq \sqrt{(x^\top V_t \hparam)^2 - \|x\|_{V_t}^2 \phi_t^2 } 
\$
Note that 
\$
(x^\top V_t \hparam_t -  \|x\|_{V_t}\phi_t ) ^2  &=  (x^\top V_t \hparam)^2  -2x^\top V_t \hparam \phi_t \|x\|_{V_t} + \|x\|_{V_t}^2 \phi_t^2 \\
&\leq (x^\top V_t \hparam)^2  - 2 \|x\|_{V_t}^2 \phi_t^2 + \|x\|_{V_t}^2 \phi_t^2\\
& = (x^\top V_t \hparam)^2  -  \|x\|_{V_t}^2 \phi_t^2
\$
where we have used the fact that $x^\top V_t \hparam > \|x\|_{V_t}\phi_t$ in the ineuqlity.  Taking square root for both sides yields the desired result, and hence $ \|t_1 x\|_{V_t}  < \phi_t$. Similarly, one can show $\phi_t  < \|t_2 x\|_{V_t}$. This concludes the proof. 


	

%\end{proof}
\endproof


\subsection{Proof of Lemma~\ref{lem:Vinv_norm}}



%\begin{proof}
\proof{}
	Let $\{\lambda_1,\lambda_2,\dots, \lambda_{N_V}\}$ be the set of distinct eigenvalues of $V$.  By the spectral theorem, $V$ can be decomposed as $V = U \Lambda U^\top  $, where the columns of $U$ consist of orthonormal eigenvectors of $V$, denoted by $\{u_{11},\dots,u_{1 n_1},u_{21}, \dots,u_{2n_2},\dots,u_{N_V 1},\dots,u_{N_V n_{N_V}}\}$, where $n_1,\dots, n_{N_V}$ are the  algebraic multiplicity of the eigenvalues respectively. Since $V$ is a symmetric matrix, the eigenvectors for a basis of $\RR^d$ and we have  $\sum_{i = 1}^{N_V}n_i = d$. 	  We can write $x$ as a linear combination 
	$
	x = \sum_{i=1}^{N_V} \sum_{j=1}^{n_i} w_{ij}  u_{ij}
	$,
	where
	$
	\sum_{i=1}^{N_V} \sum_{j=1}^{n_i} w_{ij}^2 = 1
	$. Define $a_i \coloneqq \sum_{j=1}^{n_i} w_{ij}^2 $, by direct computation, we have $ \sum_{i=1}^{N_V} a_i = 1 $  and $\|x\|_V^2 = \sum_{i=1}^{N_V} \lambda_ia_i $ and  $\|x\|_{V^{-1}}^2 = \sum_{i=1}^{N_V} \lambda_i^{-1} a_i $ .

	Now we study the range of $\|x\|_{V^{-1}}^2$ when $\|x\|_{V}^2$ is bounded as $m \leq \|x\|_{V}^2 \leq M$. First, let's focus on  maximizing $\|x\|_{V^{-1}}^2$, it suffices to solve the LP problem
	\$
	\text{maximize}& \sum_{i=1}^{N_V} a_i\lambda_i^{-1} \quad
	\text{s.t.} \quad \forall i, a_i\geq 0, \ 
	\sum_{i=1}^{N_V}  a_i = 1, \ 
	\sum_{i=1}^{N_V} a_i \lambda_i \geq m,\ 
	\sum_{i=1}^{N_V}  a_i \lambda_i \leq M.
	\$
	The Lagrangian is given by
	\$
	L = \sum_{i=1}^{N_V}  a_i\lambda_i^{-1} + \mu ( 1- \sum_{i=1}^{N_V}  a_i) + \eta (\sum_{i=1}^{N_V}  a_i \lambda_i - m)+ \gamma ( M - \sum_{i=1}^{N_V}  a_i \lambda_i )+\sum_{i=1}^{N_V}  \kappa_ia_i .
	\$
	The KKT conditions are given by
	\#\begin{cases}\label{eq:KKT}
		 & \nabla_{a_i} L  = \lambda_i^{-1} - \mu + \eta \lambda_i  -  \gamma \lambda_i+ \kappa_i =0,                         \\
		 & a_i\geq 0, \forall i, \sum_{i=1}^{N_V}  a_i = 1, \sum_{i=1}^{N_V} a_i \lambda_i \geq m, \sum_{i=1}^{N_V}  a_i \lambda_i \leq M,  \\
		 & \eta \geq  0,\gamma \geq 0,\kappa_i \geq 0, \forall i,                                                             \\
		 & \eta(\sum_{i=1}^{N_V} a_i \lambda_i - m) = 0, \gamma(M - \sum_{i=1}^{N_V} a_i \lambda_i) = 0,\kappa_i a_i = 0,  \forall i.
	\end{cases}
	\#

	To satisfy the first condition above, $\kappa_i$'s can only be zero for at most two indices.  Hence, $a_i$ can only be non-zero for at most two distinct eigenvalues, denoted by $\lambda_i$ and $\lambda_j$, where  $i<j$ and  $\lambda_i >\lambda_j$. Namely, the solution to  \eqref{eq:KKT}  lies in the subspace spanned by the eigenvectors corresponding to $\lambda_i$ and $\lambda_j$.    
	
	
	
	Let $y = \|x\|_V^2 $, we have $a_i\lambda_i + a_j \lambda_j =y$ and $a_i\lambda_i^{-1} + a_j \lambda_j ^{-1} = \|x\|_{V^{-1}}^2$. Note that $\sum_{i=1}^{N_V} a_i = a_i + a_j  = 1$, by direct computation, the closed form of $\|x\|_{V^{-1}}^2 $ is given by
	$
	\|x\|_{V^{-1}}^2 = \frac{1}{ \lambda_{i}} + \frac{1}{ \lambda_{j}}   - \frac{y}{\lambda_{i}\lambda_{j} }\eqqcolon f(y,\lambda_{i},\lambda_{j} )
	$. 
	
	Clearly, we have
	\$
	\begin{cases}
		\frac{\partial f}{\partial y}           & = -\frac{1}{\lambda_{i}\lambda_{j} } < 0,              \\
		\frac{\partial f}{\partial \lambda_{i}} & = (\frac{y}{\lambda_j} - 1)\frac{1}{\lambda_i^2 } > 0, \\
		\frac{\partial f}{\partial \lambda_{j}} & = (\frac{y}{\lambda_i} - 1)\frac{1}{\lambda_j^2 } < 0.
	\end{cases}
	\$


	Then the maximum of $\|x\|_{V^{-1}}^2 $ is obtained when $\lambda_i = \lambda_1$, $\lambda_j = \lambda_{N_V}$, $y = m$. Therefore, the solution to the LP problem is any unit vector $x^{\star}_{\max}$ that lies in the subspace spanned by the eigenvectors corresponding to $\lambda_1$ and $\lambda_{N_V}$. Moreover, we have
	\$
	\|x^\star_{\max}\|_{V^{-1}}^2 = \frac{1}{ \lambda_{1}} + \frac{1}{ \lambda_{N_V}}   - \frac{m}{\lambda_{1}\lambda_{N_V} }.
	\$

	Similarly, by investigating the KKT conditions for the LP problem that minimize $\sum_{i=1}^{N_V} a_i\lambda_i^{-1}$,   the minimum of $\|x\|_{V^{-1}}^2 $ is obtained when  $\lambda_i = \lambda_{k} $, $\lambda_j = \lambda_{k+1}$, $y = M$, where $k$ is such that $\lambda_{k} = \max_{i \in[N_v]} \{\lambda_i  \geq M\}$, and hence $\lambda_{k+1} = \min_{i \in[N_v]} \{\lambda_i  < M\}$. The solution vector is  any unit vector $x^{\star}_{\min}$ that lies in the subspace spanned by the eigenvectors corresponding to $\lambda_{k}$ and $\lambda_{k+1}$, and we have
	\$
	\|x^\star_{\min}\|_{V^{-1}}^2 = \frac{1}{ \lambda_{k}} + \frac{1}{ \lambda_{k+1} }   - \frac{M}{\lambda_{k}\lambda_{k+1}  }.
	\$
	This concludes the proof.

\endproof





\section{Other Proofs}
\label{sec:other_proofs}
\subsection{Proof of Proposition~\ref{prop:concentration} }
\label{proof:concentration}
\proof{}
	\label{pf:prop_rls}
	By Propostion~\ref{prop:beta}, we have
	\$
	\prob{\hat{\cA}_T}
	& = \prob{\cap_{t = 1}^T \left\{ \|\hparam_t - \param_t\|_{V_t}\leq \beta^{\RLS}_t(\delta^\prime) \right\}
	}  \\
	&\geq 1 - \sum_{t = 1}^T \prob{ \|\hparam_t - \param_t\|_{V_t}\leq \beta^{\RLS}_t(\delta^\prime) }\\
	&\geq 1 - \frac{\delta}{2}.
	\$
	
	Similarly, by Definition~\ref{def:tilde_beta}, we have 
		\$
	\prob{\tilde{\cA}_T}
	& = \prob{\cap_{t = 1}^T \left\{ \tilde{\theta}_t \in \cE_t^{\PVT}(\delta ^\prime) |\mathcal{F}_t\right\}
	}  \\
	&\geq 1 - \sum_{t = 1}^T \prob{ \tilde{\theta}_t \in \cE_t^{\PVT}(\delta ^\prime) |\mathcal{F}_t }\\
	&\geq 1 - \frac{\delta}{2}.
	\$
Combining the two inequalities above, we have $\prob{\cA_T} \geq 1 - \delta$. 
\endproof

\subsection{Proof of Proposition~\ref{prop:instantaneous} }
\label{proof:instantaneous}
\proof{}Since ${\theta}^\star \in \cE_t^{RLS}(\delta^\prime) $ and $\tilde{\theta}_t \in \cE_t^{\PVT}(\delta ^\prime)$, it holds that 
\#\label{eq:two_radius}
 \|{\theta}^\star - \hparam_t\|_{V_t}\leq \beta_t^{RLS}(\delta^\prime),  \quad \|\tilde{\theta}_t - \hparam_t\|_{V_t}\leq \beta^{\PVT}_t (\delta^\prime)= \iota_t \beta^{RLS}_t(\delta^\prime).
\#

In the following, we drop the dependence on $\delta^\prime$ for notation simplicity. We have 
\$
 \dotp{x_t^\star}{ \theta^{\star}} - \dotp{\tilde{x}_t}{ \theta^{\star}}
& = \left(\dotp{x^{\star}_t}{ \theta^{\star}} - \dotp{x^{\star}_t}{ \tilde{\theta}_t}\right) + \left(\dotp{x^{\star}_t}{ \tilde{\theta}_t} - \dotp{\tilde{x}_t}{  \tilde{\theta}_t}\right)\\
& \quad + \left(\dotp{\tilde{x}_t}{\tilde{\theta}_t}-\dotp{\tilde{x}_t}{ \hparam_t}\right) +  \left(\dotp{\tilde{x}_t}{\hparam_t}-\dotp{\tilde{x}_t}{ \theta^{\star}}\right).
\$


To bound the second term on the right hand side, recall by Equation \eqref{eq:POFUL} the POFUL action $\tilde{x}_t$ satisfies 
\$
\dotp{x^{\star}_t}{ \tilde{\theta}_t} +  \tau_t  \|x^{\star}_t\|_{V^{-1}_t}\beta_t^{RLS} \leq  \dotp{\tilde{x}_t}{  \tilde{\theta}_t} + \tau_t  \|\tilde{x}_t\|_{V^{-1}_t}\beta_t^{RLS}.
\$
Rearranging the ineuqlity, we obtain $\dotp{x^{\star}_t}{ \tilde{\theta}_t} - \dotp{\tilde{x}_t}{  \tilde{\theta}_t} \leq \tau_t\|\tilde{x}_t\|_{V^{-1}_t} \beta_t^{RLS}-\tau_t \|x^{\star}_t\|_{V^{-1}_t} \beta_t^{RLS}$.


The other three terms are bounded similarly using the Cauchy-Schwarz inequality,  the triangle ineuqlity of the $V_t^{-1}$-norm,  and the concentration condition \eqref{eq:two_radius}. As a result, we have
\$
\dotp{x^{\star}_t}{ \theta^{\star}} - \dotp{x^{\star}_t}{ \tilde{\theta}_t} &\leq (1+\iota_t)\|x^{\star}_t\|_{V^{-1}_t}\beta_t^{RLS} , \\ \dotp{\tilde{x}_t}{\tilde{\theta}_t}-\dotp{\tilde{x}_t}{ \hparam_t}  &\leq  \iota_t \|\tilde{x}_t\|_{V^{-1}_t}\beta^{RLS}_t,\\
\dotp{\tilde{x}_t}{\hparam_t}-\dotp{\tilde{x}_t}{ \theta^{\star}} &\leq  \|\tilde{x}_t\|_{V^{-1}_t}\beta^{RLS}_t.
\$
Combing all terms above, we have 
\$
\dotp{x_t^\star}{ \theta^{\star}} - \dotp{\tilde{x}_t}{ \theta^{\star}}
 \leq  (1+\iota_t - \tau_t)\|x^{\star}_t\|_{V^{-1}_t}\beta_t^{RLS} + (1+\iota_t + \tau_t) \|\tilde{x}_t\|_{V^{-1}_t}\beta^{RLS}_t.
\$
\endproof


\subsection{Proof of Theorem~\ref{thm:oracle_regret}}
\label{sec:proof_regret}
\proof{}
We formally prove Theorem~\ref{thm:oracle_regret} for completeness. The proof techniques are developed in previous papers \citep{abbasi2011improved,agrawal2013thompson,abeille2017linear}.

Throughout the proof, we condition on the event $\cA_T$, which holds with probability $1-\delta$ by Proposition~\ref{prop:concentration}. Applying Proposition~\ref{prop:instantaneous}, we obtain
\$
\cR(T) 
& \leq \sum_{t = 1}^T \left( \dotp{x_t^\star}{\theta^{\star}}-\dotp{\tilde{x}_t}{\theta^{\star}} \right)\event{\cA_t}\\
& \leq  \sum_{t = 1}^T  (1+\iota_t-\tau_t)\|x^{\star}_t\|_{V^{-1}_t}\beta_t^{RLS}(\delta^\prime) +   (1+\iota_t+\tau_t)\|\tilde{x}_t\|_{V^{-1}_t}\beta_t^{RLS}(\delta^\prime).
\$

Recall $\|x^{\star}_t\|_{V^{-1}_t} = \alpha_t \|\tilde{x}_t\|_{V^{-1}_t}$ and $\mu_t  =  \alpha_t(1+\iota_t -\tau_t)+1+\iota_t+\tau_t$, we have 
\$
\cR(T)  \leq \sum_{t = 1}^T   \mu_t \|\tilde{x}_t\|_{V^{-1}_t}\beta_t^{RLS}(\delta^\prime).
\$ 

Applying the Cauchy-Schwarz inequality and Proposition~\ref{prop:potential}, note that $\max_{t\in[T]}  \beta_t^{RLS}(\delta^\prime) = \beta_T^{RLS}(\delta^\prime)$, we obtain 
\$
\cR(T) &\leq   \sqrt{ \sum_{t=1}^T \|\tilde{x}_t\|_{V_t^{-1}} ^2 } \sqrt{\sum_{t=1}^T\mu_t^2 \left(\beta_t^{RLS}(\delta^\prime)\right)^2} \leq \sqrt{2 d \log \left(1+\frac{T}{\lambda}\right)}   \sqrt{\sum_{t=1}^T\mu_t^2 } \beta_T^{RLS}(\delta^\prime).
\$


This concludes the proof.
\endproof

%\paragraph{Bounding $\cR^{RLS}(T)$ and $\cR^{\PVT}(T)$.}
%Conditioned on $\cA_T$, we have
%\#\label{eq:regret_rls_ca}
%\cR^{RLS}(T) +\cR^{\PVT}(T) &\leq \sum_{t=1}^T\left(\dotp{\tilde{x}_t}{\tilde{\theta}_t}-\dotp{\tilde{x}_t}{ \hat{\theta}_t} \right) \event{\cA_t} + \left(\dotp{\tilde{x}_t}{\hat{\theta}_t}-\dotp{\tilde{x}_t}{ \theta^{\star}}\right) \event{\cA_t}  \nonumber\\
%&\leq \sum_{t=1}^T \|\tilde{x}_t\|_{V_t^{-1}} \|\theta^{\star} - \hat{\theta}_t\|_{V_t}+ \|\tilde{x}_t\|_{V_t^{-1}} \|\hat{\theta}_t - \tilde{\theta}_t\|_{V_t} \nonumber\\
%& \leq \sum_{t=1}^T \|\tilde{x}_t\|_{V_t^{-1}}  \left( \beta^{\RLS}_t(\delta^\prime)  +  \beta^{\PVT}_t(\delta^\prime)  \right)\nonumber\\
%& =  \sum_{t=1}^T \|\tilde{x}_t\|_{V_t^{-1}}  (1 + \iota_t)\beta^{RLS}_t(\delta^\prime) \#
%
%
%
%
%
%
%\paragraph{Bounding $\cR^{PE}(T)$.} 
%By direct computation,
% \$
%\left(\dotp{x^{\star}_t}{ \theta^{\star}} - \dotp{\tilde{x}_t}{  \tilde{\theta}_t} \right)\event{\cA_t}&\leq \left(\dotp{x^{\star}_t}{ \theta^{\star}} - \dotp{x^{\star}_t}{ \tilde{\theta}_t}\right)\event{\cA_t} + \left(\dotp{x^{\star}_t}{ \tilde{\theta}_t} - \dotp{\tilde{x}_t}{  \tilde{\theta}_t}\right)\event{\cA_t} \\
%%&\leq \|x^{\star}_t\|_{V^{-1}_t}(1+\iota_t)\beta_t^{RLS} +  \left(\dotp{x^{\star}_t}{ \tilde{\theta}_t} - \dotp{\tilde{x}_t}{  \tilde{\theta}_t}\right)\event{\cA_t}\\
%&\leq\|x^{\star}_t\|_{V^{-1}_t}(1+\iota_t)\beta_t^{RLS} + \|\tilde{x}_t\|_{V^{-1}_t} \tau_t\beta_t^{RLS}- \|x^{\star}_t\|_{V^{-1}_t} \tau_t\beta_t^{RLS}\\
%&\leq(1+\iota_t-\tau_t)\|x^{\star}_t\|_{V^{-1}_t}\beta_t^{RLS} +   \tau_t\|\tilde{x}_t\|_{V^{-1}_t}\beta_t^{RLS}.
%\$
%where the second inequality follows from the Cauchy-Schwarz inequality and the definition of POFUL \eqref{eq:POFUL}. We have 
%\#\label{eq:regret_pe}
%\cR^{PE}(T) &\leq \sum_{t=1}^T (1+\iota_t-\tau_t)\|x^{\star}_t\|_{V^{-1}_t}\beta_t^{RLS} +   \tau_t\|\tilde{x}_t\|_{V^{-1}_t}\beta_t^{RLS}\nonumber\\
%&\leq \sum_{t=1}^T (1+\iota_t-\tau_t)\alpha_t \|\tilde{x}_t\|_{V^{-1}_t}\beta_t^{RLS} +   \tau_t\|\tilde{x}_t\|_{V^{-1}_t}\beta_t^{RLS}.
%\#
%
%Combining  \eqref{eq:regret_rls_ca} and \eqref{eq:regret_pe}, we have
%\$
%\cR(T) &\leq  \sum_{t=1}^T \|\tilde{x}_t\|_{V_t^{-1}}  (1 + \iota_t)\beta^{RLS}_t(\delta^\prime) + \sum_{t=1}^T  (1+\iota_t-\tau_t)\alpha_t \|\tilde{x}_t\|_{V^{-1}_t}\beta_t^{RLS}(\delta^\prime) +   \tau_t\|\tilde{x}_t\|_{V^{-1}_t}\beta_t^{RLS}(\delta^\prime)\\
%& = \sum_{t=1}^T \left( (1+\iota_t-\tau_t)\alpha_t + (1+\iota_t+\tau_t)\right) \|\tilde{x}_t\|_{V^{-1}_t}\beta_t^{RLS}(\delta^\prime) \\
%& = \sum_{t=1}^T \mu_t \|\tilde{x}_t\|_{V^{-1}_t}\beta_t^{RLS}(\delta^\prime) \\
%\$
%
%Applying the Cauchy-Schwarz inequality and Proposition~\ref{prop:potential}, note that $\max_{t\in[T]}  \beta_t^{RLS}(\delta^\prime) = \beta_T^{RLS}(\delta^\prime)$, we have 
%\$
%\cR(T) &\leq   \sqrt{ \sum_{t=1}^T \|\tilde{x}_t\|_{V_t^{-1}} ^2 } \sqrt{\sum_{t=1}^T\mu_t^2 \left(\beta_t^{RLS}(\delta^\prime)\right)^2} \leq \sqrt{2 d \log \left(1+\frac{T}{\lambda}\right)}   \sqrt{\sum_{t=1}^T\mu_t^2 } \beta_T^{RLS}(\delta^\prime).
%\$
%
%
%This concludes the proof.
\endproof



\subsection{Proof of Proposition~\ref{prop:align}}
%\begin{proof}
\label{proof:align}
\proof{}
	This corresponds to the case where $m, M \to \lambda_1$ in Lemma~\ref{lem:Vinv_norm}.
	Without loss of generality, we let $M>\lambda_2$. Then
	$
	\hat{\alpha}_t  =(\lambda_{1}^{-1} + \lambda_{d}^{-1}   - m{\lambda_{1}^{-1}\lambda_{d}^{-1} })/(\lambda_{k}^{-1} +  \lambda_{k+1}^{-1}   - {M}{\lambda_{k}^{-1}\lambda_{k+1}^{-1} } ) \to \lambda_{1}^{-1}/\lambda_{1}^{-1} = 1
	$. 
%	\end{proof}
\endproof







\section{Empirical Validation of Case 2 from Section~\ref{eg:principle_eigen}}
\label{sec:vali}

%To empirically validate that Case 2 in Section~\ref{eg:principle_eigen} isn't vacant, we check the condition presented in Proposition \ref{prop:align}. We track the ratio 
%$
%\zeta_t\coloneqq \frac{\|\hat{\theta}_t\|_{V_t}^2/ \|\hat{\theta}_t\|^2}{\lambda_1(V_t)}
%$. Note that when $\hat{\theta}_t$ lies around the top eigenspace of $V_t$, we have $\|\hat{\theta}_t\|_{V_t}^2 \approx \lambda_1 \|\hat{\theta}_t\|^2$. Hence the ratio is a proxy that measures how close $\hat{\theta}_t$ is aligned with the top eigenspace of $V_t$. Specifically, $\zeta_t =1$ implies that $\hat{\theta}_t$ lies in the top eigenspace. 
%
%We rerun Example 1 in Section \ref{sec:simulation} as it represents the general linear bandit problem without special structure. 
%We plot the empirical $\{\zeta_t\}$ sequence in Figure~\ref{fig:proj}.


To empirically confirm that Case 2, as discussed in Section~\ref{eg:principle_eigen}, is not merely theoretical, we examine the condition outlined in Proposition \ref{prop:align}. Our focus is on the ratio
$
\zeta_t\coloneqq \frac{\|\hat{\theta}_t\|_{V_t}^2/ \|\hat{\theta}_t\|^2}{\lambda_1(V_t)}
$.  Notably, $\|\hat{\theta}_t\|_{V_t}^2 $ tends to approximate $\lambda_1 \|\hat{\theta}_t\|^2$ when $\hat{\theta}_t$ is proximate to the top eigenspace of $V_t$. Consequently, this ratio serves as a proxy of the alignment of $\hat{\theta}_t$ with the top eigenspace of $V_t$. Specifically, a $\zeta_t$ value of 1 signifies that $\hat{\theta}_t$ is within the top eigenspace.

We re-evaluate Example 1 from Section \ref{sec:simulation}, which epitomizes the general linear bandit problem in its standard form. The empirical sequence ${\zeta_t}$ is depicted in Figure~\ref{fig:proj}.

 \begin{figure}[htp]
	\centering
%	\subfigure{
		\includegraphics[width=.5\textwidth]{figures/proj_ex1.jpg} 
	%}
	\caption{Evolution of the alignment proxy $\zeta_t$ in Example 1.}
	\label{fig:proj}
\end{figure}

%As we see in Figure~\ref{fig:proj}, for all bandit algorithms, $\zeta_t$ is already very high (above 0.9) at the very early stage of the implementation and tends to 1 eventually. This validates the tendency described in Case 2. Moreover, Figure~\ref{fig:proj} shows that this happens for all bandit algorithms considered, coinciding with our intuition that the tendency might come from the online nature of bandit algorithms. 
%(We remark that the sharp decrease at the beginning is due to the behavior of the regularized least square estimator in over-parameterized regimes when $t$ is smaller than the dimension $d$.)



Figure~\ref{fig:proj} reveals that for every bandit algorithm, the value of $\zeta_t$ is notably high (exceeding 0.9) at the early stages of implementation and gradually converges towards 1. This observation validates the behavior outlined in Case 2. Furthermore, the consistency of this phenomenon across all examined bandit algorithms suggests that such a tendency is likely a characteristic of the online nature of these algorithms.

It is important to note that the initial sharp decline observed is attributed to the behavior of the regularized least squares estimator in over-parameterized scenarios, when $t$ is less than the dimension $d$.









